%----------------------------------------------------------------------------------------
% Tailored CV — PhD Position, Probabilistic Metabolic Flux Analysis
% HDS-LEE Graduate School, Forschungszentrum Jülich — Ref. 2025D-179
%----------------------------------------------------------------------------------------

\documentclass[10pt,a4paper,sans]{moderncv}

\usepackage[utf8]{inputenc}
\usepackage[T1]{fontenc}

\moderncvstyle{classic}
\moderncvcolor{blue}

\usepackage[scale=0.78]{geometry}
\usepackage{ragged2e}

%----------------------------------------------------------------------------------------

\firstname{Kundai Farai}
\familyname{Sachikonye}

\address{Auf der Lichtung 19}{82178 Puchheim, Germany}
\mobile{(+49) 1626386344}
\email{kundai.sachikonye@bitspark.com}
\homepage{github.com/fullscreen-triangle}{github.com/fullscreen-triangle}

%----------------------------------------------------------------------------------------

\begin{document}

\makecvtitle

%----------------------------------------------------------------------------------------
%	EDUCATION
%----------------------------------------------------------------------------------------

\section{Education}

\cventry{2019--2021}{PhD candidate, Computational Lipidomics}{Technische Universität München / Bayerisches Forschungsinstitut}{}{}{
  MS-based lipidomics at HPC scale: isotopologue quantification, multi-omics
  statistical modelling, federated inference across distributed clinical sites.
  Departure to pursue independent computational research.
}
\cventry{2018--2019}{MSc Computational Biology}{Universität Tübingen}{}{}{
  Probabilistic graphical models, statistical inference, Bayesian methods,
  computational systems biology, metabolic network analysis.
}
\cventry{2014--2017}{MSc Pharmaceutical Biotechnology}{HAW Hamburg}{}{}{}
\cventry{2011--2014}{BSc Biochemistry and Cell Biology}{Jacobs University Bremen}{}{}{}

%----------------------------------------------------------------------------------------
%	RELEVANT MANUSCRIPTS
%----------------------------------------------------------------------------------------

\section{Relevant Manuscripts}

\cvitem{2025}{
  \textbf{Backward Training Principle.}
  Formal proof: a model calibrated at the worst-case misspecification instance
  is calibrated at all less-misspecified instances by feasibility inclusion.
  Forward Training Collapse Theorem (positive VC failure gap for easy-to-hard
  direction); Sample Complexity Theorem ($(N{+}1)\times$ savings for depth-$N$
  cascades). Directly applicable to hierarchical Bayesian MFA misspecification.
  \textit{Manuscript.}
}

\cvitem{2025}{
  \textbf{On Disease Epistemology in Blind Receiver.}
  Hierarchical Bayesian model with loop-holonomy self-consistency diagnostic;
  AUC$=1.00$; sparse $\ell_1$ LP for minimum-intervention parameter selection.
  \textit{Manuscript. AIMe Registry / TUM Weihenstephan.}
}

%----------------------------------------------------------------------------------------
%	RESEARCH SOFTWARE
%----------------------------------------------------------------------------------------

\section{Research Software}

\cvitem{purpose + homo-veloce}{
  \textbf{Backward Training Principle --- Misspecification-Robust Model Training.}
  Formal proof that calibration at the hardest (most misspecified) instance
  propagates to all less-misspecified sub-models by feasibility inclusion.
  $\mathcal{C}_\mathrm{hard}\subseteq\mathcal{C}_\mathrm{easy}$ implies
  $\epsilon$-competence transfers downward without retraining.
  Forward Training Collapse Theorem: positive VC failure gap for the conventional
  easy-to-hard calibration direction.
  homo-veloce: distils Gaussian process and physics-based solvers into lightweight
  neural surrogates via knowledge distillation --- same pattern as neural ODE
  surrogates for metabolic network simulation.
  \href{https://github.com/fullscreen-triangle/purpose}{github.com/fullscreen-triangle/purpose}
}

\cvitem{gospel}{
  \textbf{Hierarchical Bayesian Multi-Omics Inference.}
  Bayesian network inference for multi-omics integration; variational Bayes
  posterior approximation at $\mathcal{O}(10^6)$ latent variable scale;
  Gaussian mixture models for population structure; continuous shrinkage priors.
  Validated: 1000~Genomes, gnomAD~v3.1.2, UK~Biobank.
  Architecture applicable to hierarchical flux posterior inference with
  stoichiometric network priors.
  \href{https://github.com/fullscreen-triangle/gospel}{github.com/fullscreen-triangle/gospel}
}

\cvitem{syndrome}{
  \textbf{Probabilistic State Inference with Sparse Regularisation.}
  Eight-dimensional latent state $\mathbf{D}\in[0,1]^8$ with coherence indices
  $\eta_i\in[0,1]$; Bayesian P(R\,|\,Q) posterior gating; $\mathcal{O}(N\log N)$
  trajectory; sparse $\ell_1$ LP for minimum-norm parameter selection.
  AUC$=1.00$ in 25/25 validation experiments.
  Analogous to flux vector inference under thermodynamic feasibility constraints.
  \href{https://github.com/fullscreen-triangle/syndrome}{github.com/fullscreen-triangle/syndrome}
}

\cvitem{lavoisier}{
  \textbf{Metabolomics / Lipidomics Data Pipeline (MFA Input Layer).}
  Dual-pipeline MS analysis: numerical isotopologue quantification + CNN
  spectral classification (PyTorch); 98.9\,\% NIST accuracy; Ray-distributed
  HPC; $>$100\,GB mzML support; memory-mapped I/O.
  Processes the primary measurement data for $^{13}$C-MFA:
  isotopically labelled mass spectra at production scale.
  \href{https://github.com/fullscreen-triangle/lavoisier}{github.com/fullscreen-triangle/lavoisier}
}

\cvitem{four-sided-triangle}{
  \textbf{HPC-Scale Inference Pipeline.}
  Rust core (15--50$\times$ speedup via PyO3 bindings); Ray-distributed
  parallel chain execution; FastAPI orchestration; Docker/Kubernetes.
  Demonstrates the Python + Rust + Ray HPC architecture applicable to
  large-scale MCMC and variational inference for metabolic flux posteriors.
  \href{https://github.com/fullscreen-triangle/four-sided-triangle}{github.com/fullscreen-triangle/four-sided-triangle}
}

%----------------------------------------------------------------------------------------
%	EXPERIENCE
%----------------------------------------------------------------------------------------

\section{Experience}

\cventry{2021--present}{Independent AI \& Computational Biology Developer}{Bitspark GmbH}{}{}{
  Hierarchical Bayesian modelling, differentiable probabilistic inference,
  HPC-scale distributed sampling, metabolomics and genomics data pipelines.
  All systems publicly deployed and validated.
}

\cventry{2019--2021}{Computational Lipidomics}{TUM / Bayerisches Forschungsinstitut}{Munich}{}{
  MS-based lipidomics: isotopologue quantification, peak detection, ion annotation
  (LIPID MAPS, LipidBlast), multi-omics statistical modelling, federated
  inference for distributed clinical data. Python, Java.
}

\cventry{2019}{Process Optimisation}{Fraunhofer Institut IGB}{Stuttgart}{}{
  Real-time sensor data acquisition; dynamic programming and LP for
  photobioreactor optimisation. Python, C.
}

%----------------------------------------------------------------------------------------
%	TECHNICAL SKILLS
%----------------------------------------------------------------------------------------

\section{Technical Skills}

\cvitem{Bayesian Methods}{
  Variational Bayes (ELBO optimisation); MCMC (NUTS, HMC, Metropolis-Hastings);
  hierarchical / multi-level models; continuous shrinkage priors; probabilistic
  graphical models; Bayesian network inference; automatic differentiation
  (PyTorch \texttt{autograd}); normalising flows; Laplace approximation
}

\cvitem{Metabolic Flux Analysis}{
  Stoichiometric network modelling; isotopologue balance equations;
  $^{13}$C-MFA data structures (mzML, isotopologue distributions);
  matrix effects and measurement uncertainty in MS-based quantification;
  LIPID MAPS, LipidBlast, HMDB metabolite databases
}

\cvitem{HPC / Scientific Computing}{
  Python (primary); Rust with PyO3 bindings (15--50$\times$ speedup);
  C (Fraunhofer context); Ray distributed sampling; Dask; Docker, Kubernetes;
  memory-mapped I/O; Linux HPC environment; SLURM-compatible pipelines
}

\cvitem{ML / Probabilistic Programming}{
  PyTorch; JAX-compatible automatic differentiation; scikit-learn;
  knowledge distillation from physics / ODE solvers to neural surrogates;
  sparse $\ell_1$ regularisation; Gaussian processes
}

%----------------------------------------------------------------------------------------
%	LANGUAGES
%----------------------------------------------------------------------------------------

\section{Languages}

\cvitemwithcomment{English}{Native / Fluent}{}
\cvitemwithcomment{German}{Conversational}{Living in Germany since 2011}

%----------------------------------------------------------------------------------------

\renewcommand{\listitemsymbol}{-~}

\end{document}
